% GENERATED FILE. Edit canonical lessons, then run npm run generate:books.
% canonical-source-hash: fnv1a64:c6af03ee48734f33
% canonical-lessons: AR-C288-dhahab-wa-iyab, AR-C288-taira, AR-W288-hamza-on-yaa, AR-C288-masul, AR-C288-fatani, AR-R288-at-the-station

\chapter{At The Station}
\label{ch:ar-at-the-station}
\hlchaptermodality{288}

\begin{chapteropening}
\textbf{By the end of this chapter:} \emph{I can buy a return ticket, report a missed train or a lost bag, and ask for the person in charge.}

Complete AR-R288-at-the-station: buy a return ticket, report a missed train, and ask for the person in charge.
\end{chapteropening}

\section[\texorpdfstring{\ar{تذكرة} \ar{ذهاب} \ar{وإياب} (tadhkirat dhahāb …)}{تذكرة ذهاب وإياب (tadhkirat dhahāb …)}]{\ar{تذكرة} \ar{ذهاب} \ar{وإياب} (tadhkirat dhahāb wa-iyāb) — a return ticket}

\label{lesson:AR-C288-dhahab-wa-iyab}



\begin{quote}
Give the way to the station in two orders. (\emph{Idhhab mustaqīman, thumma khudh ash-shāriʿa al-awwala ʿalā al-yamīn}.)

You are at the station now, at the counter, and the person behind the glass is waiting for you to speak.
\end{quote}

\subsection*{You'll want to know: \ar{تذكرة} \ar{ذهاب} \ar{وإياب}}
\textbf{\ar{تذكرة} \ar{ذهاب} \ar{وإياب}} — \emph{tadhkirat dhahāb wa-iyāb} — \textquotedblleft{}a return ticket\textquotedblright{}, word for word \textquotedblleft{}a ticket of going and coming back\textquotedblright{}.

\begin{quote}
\textbf{\ar{أريد} \ar{تذكرة} \ar{ذهاب} \ar{وإياب،} \ar{من} \ar{فضلك}.} — \emph{urīdu tadhkirat dhahāb wa-iyāb, min faḍlik} — I'd like a return ticket, please.
\end{quote}

\emph{Tadhkira} ends in \textbf{\ar{ة}}, and before the noun that follows it you hear its \emph{t}: \emph{tadhkirat dhahāb}. For one way only, say \emph{tadhkirat dhahāb faqaṭ}, \textquotedblleft{}a ticket of going only\textquotedblright{}.

\begin{grammarlens}[title={Grammar Lens: three questions at the counter}]
\begin{quote}
\textbf{\ar{بكم} \ar{التذكرة؟}} — \emph{bi-kam at-tadhkira?} — How much is the ticket?
\end{quote}

\begin{quote}
\textbf{\ar{من} \ar{أي} \ar{رصيف؟}} — \emph{min ayyi raṣīf?} — From which platform?
\end{quote}

\begin{quote}
\textbf{\ar{متى} \ar{ينطلق} \ar{القطار؟}} — \emph{matā yanṭaliqu al-qiṭār?} — When does the train leave?
\end{quote}

\emph{Bi-kam} is \textquotedblleft{}for how much\textquotedblright{}, \emph{ayy} is \textquotedblleft{}which\textquotedblright{}, and \emph{matā} is \textquotedblleft{}when\textquotedblright{} in a question. \emph{Yanṭaliqu} is the present of \emph{inṭalaqa}, to set off.
\end{grammarlens}

\begin{cousinweb}
\textbf{\ar{ذهاب}} (\emph{dhahāb}) is the verbal noun of \emph{dhahaba}, \textquotedblleft{}to go\textquotedblright{}, on the root \emph{dh-h-b}. \textbf{\ar{إياب}} (\emph{iyāb}) is built on the native Arabic root \emph{ʾ-w-b}, \textquotedblleft{}to return from a journey\textquotedblright{} (Lane), who gives \emph{iyāb} among its verbal nouns and notes that a \emph{y} has taken the place of the root's \emph{w}. The pair rhymes, which is part of why it sticks.
\end{cousinweb}

\subsection*{Your turn}
Say these aloud:

\begin{itemize}
\raggedright
  \item \textquotedblleft{}I'd like a return ticket, please\textquotedblright{} — \emph{urīdu tadhkirat dhahāb wa-iyāb, min faḍlik}
  \item \textquotedblleft{}one way only\textquotedblright{} — \emph{tadhkirat dhahāb faqaṭ}
  \item \textquotedblleft{}How much is the ticket?\textquotedblright{} — \emph{bi-kam at-tadhkira?}
  \item \textquotedblleft{}From which platform?\textquotedblright{}, then \textquotedblleft{}when does the train leave?\textquotedblright{}
  \item \textquotedblleft{}I want to buy two tickets\textquotedblright{} — \emph{urīdu an ashtariya tadhkiratayn}
\end{itemize}

\begin{tcolorbox}[breakable,colback=teal!4,colframe=teal!35!black,title={Before you move on}]
\textquotedblleft{}A return ticket\textquotedblright{}? (\emph{\textbf{Tadhkirat dhahāb wa-iyāb.}}) \textquotedblleft{}One way\textquotedblright{}? (\emph{\textbf{Tadhkirat dhahāb faqaṭ.}}) \textquotedblleft{}How much?\textquotedblright{} (\emph{\textbf{Bi-kam?}}) What does \emph{iyāb} mean? (\textbf{Coming back.})
\end{tcolorbox}
\section[\texorpdfstring{\ar{طائرة} (ṭāʾira)}{طائرة (ṭāʾira)}]{\ar{طائرة} (ṭāʾira) — a plane}

\label{lesson:AR-C288-taira}



\begin{quote}
Ask for a return ticket, then ask the price. (\emph{Urīdu tadhkirat dhahāb wa-iyāb}, \emph{bi-kam at-tadhkira?})

Say the Arabic for an airport and a pilot. (\emph{Maṭār}, \emph{ṭayyār}.) They share a root, and so does the thing between them.
\end{quote}

\subsection*{You'll want to know: \ar{طائرة}}
\textbf{\ar{طائرة}} — \emph{ṭāʾira} — \textquotedblleft{}a plane, an aeroplane\textquotedblright{}.

\begin{quote}
\textbf{\ar{وصلت} \ar{الطائرة} \ar{متأخرة}.} — \emph{waṣalat aṭ-ṭāʾira mutaʾakhkhira} — The plane arrived late.
\end{quote}

\begin{quote}
\textbf{\ar{سافرنا} \ar{بالطائرة}.} — \emph{sāfarnā bi-ṭ-ṭāʾira} — We travelled by plane.
\end{quote}

The verb takes the \emph{-at} of \textquotedblleft{}she\textquotedblright{}: \emph{waṣalat}. A word ending in \textbf{\ar{ة}} is feminine, and so is the verb that goes with it.

Say it \emph{ṭā-ʾi-ra}: the long \emph{ā}, then the catch, then \emph{ira}.

\begin{cousinweb}
\textbf{\ar{طائرة}} (\emph{ṭāʾira}) is built on the native Arabic root \emph{ṭ-y-r}, \textquotedblleft{}to fly\textquotedblright{} (Lane), the root of \emph{ṭāra}, \textquotedblleft{}it flew\textquotedblright{}, \emph{maṭār}, \textquotedblleft{}an airport\textquotedblright{}, and \emph{ṭayyār}, \textquotedblleft{}a pilot\textquotedblright{}. Lane gives \emph{ṭāʾir} as a flying thing, a bird, and notes that the feminine \emph{ṭāʾira} is seldom used of a female bird; the sense \textquotedblleft{}a plane\textquotedblright{} belongs to the age of planes. The catch in the middle is a hamza, on a seat the next lesson puts in your hand.
\end{cousinweb}

\subsection*{Your turn}
Say these aloud:

\begin{itemize}
\raggedright
  \item \emph{ṭāʾira}, slowly — \emph{ṭā-ʾi-ra}
  \item \textquotedblleft{}The plane arrived late\textquotedblright{} — \emph{waṣalat aṭ-ṭāʾira mutaʾakhkhira}
  \item \textquotedblleft{}We travelled by plane\textquotedblright{} — \emph{sāfarnā bi-ṭ-ṭāʾira}
  \item the four words of the root — \emph{ṭāra, maṭār, ṭayyār, ṭāʾira}
\end{itemize}

\begin{tcolorbox}[breakable,colback=teal!4,colframe=teal!35!black,title={Before you move on}]
What does \ar{طائرة} mean? (\textquotedblleft{}A plane\textquotedblright{}.) Which verb shares its root? (\emph{\textbf{Ṭāra}}, it flew.) \textquotedblleft{}By plane\textquotedblright{}? (\emph{\textbf{Bi-ṭ-ṭāʾira.}})
\end{tcolorbox}
\section[\texorpdfstring{\ar{ئ} (hamza on yāʾ)}{ئ (hamza on yāʾ)}]{\ar{ئ} (hamza on yāʾ) — the catch inside \ar{طائرة}}

\label{lesson:AR-W288-hamza-on-yaa}



\begin{quote}
Before the new one: write \ar{ي} and \ar{ؤ} once each, and name them.

Now say \textbf{\ar{طائرة}} (\emph{ṭāʾira}), \textquotedblleft{}a plane\textquotedblright{}. You met it a moment ago. This lesson takes its catch out and puts it in your hand.
\end{quote}

\subsection*{Script you'll notice: \ar{ئ}}
\textbf{\ar{ئ}} — \emph{hamza} on a \emph{yāʾ} seat, the same catch as \textbf{\ar{ؤ}}, written \emph{ʾ} in this book's romanization.

It is the third letter of \textbf{\ar{طائرة}}, after \textbf{\ar{ا}}. Because \textbf{\ar{ا}} does not join to the letter after it, \textbf{\ar{ئ}} starts a new piece. It is two pieces itself: below, a small tooth like the tooth of \textbf{\ar{ي}} inside a word, with no dots — that is the seat; on top, the hamza mark. The tooth joins on to the \textbf{\ar{ر}} after it.

The sound is the catch you already say in \emph{suʾāl} and \emph{saʾala}. Here the \emph{i} right after it calls for the \textbf{\ar{ي}} seat.

\subsection*{Writing: \ar{ئ}}
Put your pen on \ar{ئ} and follow its line. Write the dotless tooth first, then add the hamza above it. Copy the shape you can see — slowly, and larger than it is printed.

\begin{quote}
This book does not yet tell you \textbf{where to start this character or which way to travel}. It is not written down here until it can be written down with a source. Copying what is in front of you needs no such source, and it is how the shape gets into your hand in the meantime.
\end{quote}

\subsection*{Your turn}
\begin{itemize}
\raggedright
  \item \emph{Find:} the seat and the mark of \ar{ئ} inside \textbf{\ar{طائرة}}
  \item \emph{Trace it:} \ar{ئ} three times, saying the catch as you finish each one
  \item \emph{Write it:} \ar{طا}, then a new piece: \ar{ئرة}
  \item \emph{Write it:} \textbf{\ar{طائرة}}, then \textbf{\ar{سؤال}}: one catch, two seats
  \item \emph{Read it:} \textbf{\ar{طائرة}} aloud once more
\end{itemize}

\begin{tcolorbox}[breakable,colback=teal!4,colframe=teal!35!black,title={Before you move on}]
What is this — \ar{ئ}? (\emph{\textbf{Hamza}} on a \emph{yāʾ} seat, a catch.) Which word did it come from? (\textbf{\ar{طائرة}}, \emph{ṭāʾira}, \textquotedblleft{}a plane\textquotedblright{}.) Does the seat have dots? (\textbf{No}: only the hamza on top.)
\end{tcolorbox}
\section[\texorpdfstring{\ar{مسؤول} (masʾūl)}{مسؤول (masʾūl)}]{\ar{مسؤول} (masʾūl) — the person in charge}

\label{lesson:AR-C288-masul}



\begin{quote}
Say \textquotedblleft{}a question\textquotedblright{}, then \textquotedblleft{}a plane\textquotedblright{}. (\emph{Suʾāl}, \emph{ṭāʾira}.) Which seat carries the catch in each? (\textbf{\ar{و}}, then the dotless tooth.)

When a plane is late and the employee at the desk cannot help, you ask for somebody else.
\end{quote}

\subsection*{You'll want to know: \ar{مسؤول}}
\textbf{\ar{مسؤول}} — \emph{masʾūl} — \textquotedblleft{}the person in charge, an official\textquotedblright{}; as a description, \textquotedblleft{}responsible\textquotedblright{}.

\begin{quote}
\textbf{\ar{أريد} \ar{أن} \ar{أتكلم} \ar{مع} \ar{المسؤول}.} — \emph{urīdu an atakallama maʿa al-masʾūl} — I want to speak to the person in charge.
\end{quote}

\begin{quote}
\textbf{\ar{من} \ar{المسؤول} \ar{هنا؟}} — \emph{man al-masʾūl hunā?} — Who is in charge here?
\end{quote}

As a description it takes \emph{ʿan}, \textquotedblleft{}for\textquotedblright{}: \emph{al-muwaẓẓaf masʾūl ʿan at-tadhākir}, the employee is responsible for the tickets. It is written with the \textbf{\ar{ؤ}} of \textbf{\ar{سؤال}}; some writers spell it \textbf{\ar{مسئول}}, with the other seat.

\begin{cousinweb}
\textbf{\ar{مسؤول}} (\emph{masʾūl}) is built on the native Arabic root \emph{s-a-l}, \textquotedblleft{}to ask\textquotedblright{} (Lane), the root of \emph{saʾala} and \emph{suʾāl}. Its shape \emph{mafʿūl} is the one that means \textquotedblleft{}done to\textquotedblright{}, so \emph{masʾūl} is \textquotedblleft{}asked\textquotedblright{}, and Lane lists it among the root's words for a thing asked for. \textquotedblleft{}The one asked to answer\textquotedblright{}, and so \textquotedblleft{}the one in charge\textquotedblright{}, is the usual way the modern sense is explained.
\end{cousinweb}

\subsection*{Your turn}
Say these aloud:

\begin{itemize}
\raggedright
  \item \emph{masʾūl}, slowly — \emph{mas-ʾūl}
  \item \textquotedblleft{}I want to speak to the person in charge\textquotedblright{} — \emph{urīdu an atakallama maʿa al-masʾūl}
  \item \textquotedblleft{}Who is in charge here?\textquotedblright{} — \emph{man al-masʾūl hunā?}
  \item \textquotedblleft{}Where is the person in charge of the station?\textquotedblright{} — \emph{ayna masʾūl al-maḥaṭṭa?}
\end{itemize}

\begin{tcolorbox}[breakable,colback=teal!4,colframe=teal!35!black,title={Before you move on}]
What does \ar{مسؤول} mean? (\textquotedblleft{}The person in charge; responsible\textquotedblright{}.) Which word for a question shares its root? (\textbf{\ar{سؤال}}, \emph{suʾāl}.) \textquotedblleft{}Responsible for\textquotedblright{}? (\emph{\textbf{Masʾūl ʿan.}})
\end{tcolorbox}
\section[\texorpdfstring{\ar{فاتني} \ar{القطار} (fātanī al-qiṭār)}{فاتني القطار (fātanī al-qiṭār)}]{\ar{فاتني} \ar{القطار} (fātanī al-qiṭār) — I missed the train}

\label{lesson:AR-C288-fatani}



\begin{quote}
Say \textquotedblleft{}I want to speak to the person in charge\textquotedblright{}. (\emph{Urīdu an atakallama maʿa al-masʾūl}.) Then ask for a return ticket. (\emph{Tadhkirat dhahāb wa-iyāb}.)

Now something has gone wrong. Arabic reports it in a shape English does not use.
\end{quote}

\begin{grammarlens}[title={Grammar Lens: the thing did it}]
\begin{quote}
\textbf{\ar{فاتني} \ar{القطار}.} — \emph{fātanī al-qiṭār} — I missed the train.
\end{quote}

Word for word it is \textquotedblleft{}the train passed me by\textquotedblright{}. \emph{Fāta} is \textquotedblleft{}it passed, it slipped by\textquotedblright{}, \emph{-nī} is \textquotedblleft{}me\textquotedblright{}, and the train is the subject, after its verb. Your own mistake is told as the train's doing.

\begin{quote}
\textbf{\ar{ضاعت} \ar{حقيبتي}.} — \emph{ḍāʿat ḥaqībatī} — My bag is lost.
\end{quote}

\emph{Ḍāʿa} is \textquotedblleft{}to get lost\textquotedblright{}, and the bag, feminine, gives the verb its \emph{-at}. So \emph{ḍāʿat amtiʿatī}, my luggage is lost; \emph{taʾakhkharat aṭ-ṭāʾira}, the plane is late. Each opens with the past verb and then names the thing.
\end{grammarlens}

\begin{grammarlens}[title={Grammar Lens: asking for help}]
\begin{quote}
\textbf{\ar{ماذا} \ar{يجب} \ar{أن} \ar{أعمل؟}} — \emph{mādhā yajibu an aʿmala?} — What should I do?
\end{quote}

\emph{Mādhā} is \textquotedblleft{}what\textquotedblright{}, before a verb. Open with \emph{ʿafwan}, \textquotedblleft{}excuse me\textquotedblright{}, and close, if you need to, with \emph{urīdu an atakallama maʿa al-masʾūl}.
\end{grammarlens}

\subsection*{Your turn}
Say these aloud:

\begin{itemize}
\raggedright
  \item \textquotedblleft{}I missed the train\textquotedblright{} — \emph{fātanī al-qiṭār}
  \item \textquotedblleft{}My bag is lost\textquotedblright{} — \emph{ḍāʿat ḥaqībatī}
  \item \textquotedblleft{}The plane is late\textquotedblright{} — \emph{taʾakhkharat aṭ-ṭāʾira}
  \item \textquotedblleft{}Excuse me, what should I do?\textquotedblright{} — \emph{ʿafwan, mādhā yajibu an aʿmala?}
  \item the whole report — \emph{ʿafwan, fātanī al-qiṭār, wa-ḍāʿat ḥaqībatī. mādhā yajibu an aʿmala?}
\end{itemize}

\begin{tcolorbox}[breakable,colback=teal!4,colframe=teal!35!black,title={Before you move on}]
\textquotedblleft{}I missed the train\textquotedblright{}? (\emph{\textbf{Fātanī al-qiṭār.}}) Who is the subject in it? (\textbf{The train}: it passed me by.) Why \emph{ḍāʿat} and not \emph{ḍāʿa}? (\textbf{The bag is feminine.}) \textquotedblleft{}What should I do?\textquotedblright{} (\emph{\textbf{Mādhā yajibu an aʿmala?}})
\end{tcolorbox}
\section[(review)]{Review — at the station, when things go wrong}

\label{lesson:AR-R288-at-the-station}



\begin{quote}
Ask the way to the airport. (\emph{Min faḍlik, kayfa aṣilu ilā al-maṭār?}) Then give the answer yourself: \textquotedblleft{}go straight on, then cross the bridge\textquotedblright{}. (\emph{Idhhab mustaqīman, thumma uʿbur al-jisr}.)
\end{quote}

\begin{grammarlens}[title={What you've built}]
Five lessons, and a whole trip through a station. At the counter: \emph{tadhkirat dhahāb wa-iyāb}, \emph{bi-kam?}, \emph{min ayyi raṣīf?} When it goes wrong: \emph{fātanī al-qiṭār}, \emph{ḍāʿat ḥaqībatī}, \emph{taʾakhkharat aṭ-ṭāʾira}, and \emph{mādhā yajibu an aʿmala?} And if nobody can help, \emph{al-masʾūl}.

The hamza now has all its seats in your hand: \textbf{\ar{أ}} in \textbf{\ar{سأل}}, \textbf{\ar{ؤ}} in \textbf{\ar{سؤال}} and \textbf{\ar{مسؤول}}, \textbf{\ar{ئ}} in \textbf{\ar{طائرة}}. One sound, three ways to write it.
\end{grammarlens}

\subsection*{Your turn}
\begin{itemize}
\raggedright
  \item \emph{Say it:} \textquotedblleft{}A return ticket, please. From which platform?\textquotedblright{}
  \item \emph{Say it:} \textquotedblleft{}Excuse me, I missed the plane. What should I do?\textquotedblright{} — \emph{ʿafwan, fātatnī aṭ-ṭāʾira. mādhā yajibu an aʿmala?}
  \item \emph{Say it:} \textquotedblleft{}My luggage is lost. Who is in charge here?\textquotedblright{}
  \item \emph{Say it:} \textquotedblleft{}I won't travel today; the reason is that the train is late\textquotedblright{} — \emph{lan usāfira al-yawm; as-sababu huwa anna al-qiṭāra mutaʾakhkhir}
  \item \emph{Say it:} \textquotedblleft{}I'll take the bus\textquotedblright{} — \emph{sa-ākhudhu al-ḥāfila}
  \item \emph{Say it:} \textquotedblleft{}The pilot was on a phone call\textquotedblright{} — \emph{kāna aṭ-ṭayyār fī mukālama}
  \item \emph{Write it:} \textbf{\ar{سؤال}}, \textbf{\ar{مسؤول}} and \textbf{\ar{طائرة}}, and circle each seat
\end{itemize}

\begin{tcolorbox}[breakable,colback=teal!4,colframe=teal!35!black,title={Before you move on}]
\textquotedblleft{}One way only\textquotedblright{}? (\emph{\textbf{Tadhkirat dhahāb faqaṭ.}}) \textquotedblleft{}My bag is lost\textquotedblright{}? (\emph{\textbf{Ḍāʿat ḥaqībatī.}}) \textquotedblleft{}The person in charge\textquotedblright{}? (\emph{\textbf{Al-masʾūl.}}) Which seat does the hamza take in \textbf{\ar{طائرة}}? (\textbf{\ar{ئ}}, the dotless tooth.)
\end{tcolorbox}
